% !TEX root = tarasov-qasp-zh.tex

\chapter[你对牛顿运动定律理解得如何？]{你对牛顿运动定律\\理解得如何？\label{ch-04}}

\begin{dialogue}
\tchr \hlred{请叙述牛顿第一定律。}

\std 物体保持静止状态或匀速直线运动状态，直到其他物体的作用迫使它改变这种状态为止。

\tchr \hlred{这一定律在一切参考系中都成立吗？}

\std 我不明白您的问题。

\tchr 说一个物体静止，指的是它相对于另外某个物体不动，这个物体在给定情形中就充当参考系（frame of reference）。不指明参考系而去谈论物体处于静止或某种确定的运动，是毫无意义的。物体运动的性质取决于参考系的选取。比如，行进中的车厢地板上放着一个物体：在固连于车厢的参考系里它是静止的，而在固连于铁轨的参考系里它是运动的。现在可以回到我的问题了：牛顿第一定律在一切参考系中都成立吗？

\std 呃，大概是吧。

\tchr 看来这个问题把你问住了。实验表明，牛顿第一定律并非在一切参考系中都成立。就拿车厢地板上的物体为例。我们忽略物体与地板之间的摩擦。先考察物体相对于固连于车厢的参考系的情形。可以观察到：物体本来静止在地板上，突然之间开始沿地板滑动，尽管看不出有任何作用。这里显然违背了牛顿第一定律。对这种现象的通行解释是：原来匀速直线行驶的车厢因列车制动而开始减速，而物体由于没有摩擦，相对于铁轨仍保持匀速直线运动状态。由此可以得出结论：牛顿定律在固连于铁轨的参考系中成立，而在固连于减速中的车厢的参考系中不成立。

\hlblue{牛顿第一定律在其中成立的参考系称为惯性参考系（inertial frame of reference），不成立的称为非惯性参考系（non-inertial frame of reference）。}对我们研究的大多数现象来说，可以认为固连于地面、或固连于任何相对于地面静止或做匀速直线运动的物体的参考系都是惯性系。非惯性参考系则是以加速度（或减速度）运动的参考系，例如转动的系统、加速或减速的电梯等。注意，在非惯性参考系中不成立的不仅有牛顿第一定律，还有第二定律（因为第一定律是第二定律的特例）。
\pnote{惯性参考系与非惯性参考系超出中学教学要求：中学课程默认取地面为参考系并把它当作惯性系，不展开这一概念；非惯性系与惯性力属于大学普通物理（物理竞赛常涉及）的内容。另外，由于地球自转，地面参考系其实只是近似的惯性系，但对中学和大多数普通物理问题，这一近似完全够用。}

\std 可要是牛顿定律对加速运动的参考系不适用，那在这种参考系中怎样处理力学问题呢？

\tchr 对于非惯性参考系，牛顿运动定律仍然可以使用。不过，为此必须纯粹形式地对物体再附加一个力。这个力就是所谓的惯性力（inertial force），它等于物体质量与参考系加速度的乘积，方向与物体的加速度相反。\hlblue{我要强调，这种力实际上并不存在；但只要把它形式地引入，牛顿运动定律在非惯性参考系中便照常成立。}
\pnote{原书说惯性力的方向“与物体的加速度相反”，这里的“物体”应理解为充当参考系的那个物体（参考体，如车厢、电梯），即惯性力与参考系的加速度相反，而不是与所研究物体的加速度相反——在非惯性系中，所研究的物体可以静止（加速度为零），此时惯性力并不为零。平动加速的参考系中，惯性力的大小为 $ma_0$（$a_0$ 为参考系的加速度），方向与参考系加速度的方向相反。}

我还是要建议你：解题时只用惯性参考系。这样，你要处理的各个力就都是真实存在的力。

\std 可要是我们局限于惯性参考系，那就无法分析诸如放在转动圆盘上的物体这类问题了。

\tchr 为什么无法分析？参考系由你自己选择。在这类问题中，若取固连于圆盘的参考系（即非惯性系），物体就被看作静止的；若取固连于地面的参考系（即惯性系），物体就被看作是在做圆周运动。

我仍建议你选惯性参考系。现在，\hlred{请叙述牛顿第二定律。}

\std 这个定律可以写成 $F=ma$，其中 $F$ 是作用在物体上的力，$m$ 是物体的质量，$a$ 是加速度。

\tchr 你这种简短的回答很典型。我要对你的表述提出三点批评意见：两点不大重要，一点是关键性的。第一，\hlblue{不是力由加速度产生，恰恰相反，加速度是所施力的结果。}因此，更合乎逻辑的是把这一定律的方程写成
\begin{equation}%
a=\frac{BF}{m}
\label{eq-10}
\end{equation}
其中 $B$ 是比例系数（proportionality factor），取决于式~\eqref{eq-10} 中各物理量测量单位的选择。注意，你的表述中没有提到比例系数 $B$。第二，使物体获得加速度的是作用在它上面的所有的力（尽管其中一些力可能相互平衡）。因此，\hlblue{在表述这一定律时，不应使用“力”这个词，而应使用更准确的术语“合力”（resultant force）。}
\pnote{人教版必修第一册第四章（4.3 节“牛顿第二定律”）直接给出 $F=ma$（$F$ 为合力），不出现比例系数 $B$：在国际单位制（SI）中，力的单位牛顿正是按这一定律规定的——使 $1\;\mathrm{kg}$ 物体产生 $1\;\mathrm{m/s^{2}}$ 加速度的力定义为 $1\;\mathrm{N}$，于是 $B=1$。只有各自独立地选取力、质量、加速度的单位时，才需要 $B$ 来配平；单位制问题见同册 4.4 节“力学单位制”。}

我的第三点意见最重要。牛顿第二定律确立的是力与加速度之间的关系。但力与加速度都是矢量（vector），不仅由数值（大小）表征，也由方向表征。\hlblue{你的表述没有指明方向，这是一个本质性的缺点。}你的表述丢掉了牛顿第二定律中一个至关重要的部分。正确的表述是：物体的加速度与作用在该物体上所有力的合力成正比，与物体的质量成反比，加速度的方向沿合力的方向。这一表述可以用公式写成
\begin{equation}%
\vec{a}=\frac{B\vec{F}}{m}
\label{eq-11}
\end{equation}
（字母上的箭头表示矢量）。

\std 在\chap{ch-02} 讨论作用在斜抛物体上的力时，您说过以后要说明：物体运动的方向不一定与作用在它上面的力的方向一致。您当时援引的就是牛顿第二定律吧。
\begin{figure}[!h]%
\centering
\includegraphics[width=0.4\textwidth]{figs/ch-04/fig-018a.pdf}
\caption[用矢量作图表示速度的变化。]{用矢量作图表示速度的变化。\label{fig-18}}
\end{figure}

\tchr 对，我记得；现在回到这个问题上来正合适。先回忆一下什么是加速度。我们知道，加速度由单位时间内速度的变化来表征。\figr{fig-18} 中画出了物体在两个相近时刻 $t$ 和 $t+\Delta t$ 的速度矢量 $\vec{v_{1}}$ 与 $\vec{v_{2}}$。时间 $\Delta t$ 内速度的变化是矢量 $\Delta \vec{v} =\vec{v_{2}} - \vec{v_{1}}$。按定义，加速度为
\begin{equation}%
\vec{a}(t) \cong \lim_{\Delta t \rightarrow 0}\frac{\Delta \vec{v}}{\Delta t}
\label{eq-12}
\end{equation}
或者更严格地写成
\begin{equation}%
\vec{a}(t) = \lim_{\Delta t \rightarrow 0}\frac{\Delta \vec{v}}{\Delta t}
\label{eq-13}
\end{equation}
由此可见，加速度矢量沿矢量 $\Delta v$ 的方向，后者表示在足够短的时间间隔内速度的变化。从\figr{fig-18} 明显可见，速度矢量与速度变化矢量的取向可以完全不同。这就是说，\hlblue{在一般情形下，加速度矢量与速度矢量的取向也不同。}清楚吗？

\std 是的，现在我明白了。例如，物体沿圆周运动时，速度沿圆的切线方向，而加速度沿半径指向转动中心——我说的是向心加速度（centripetal acceleration）。
\pnote{“速度沿切线、加速度指向圆心”只对匀速圆周运动成立。一般的（变速）圆周运动中，加速度除指向圆心的向心分量外，还有沿切线方向、改变速率大小的分量。学生举的正是匀速圆周运动的例子；原书此处从简，未加区分。}

\tchr 你的例子很恰当。现在回到关系式（式~\eqref{eq-11}），并弄清楚：沿作用力方向的是加速度而不是速度；与该力的大小相联系的，也仍是加速度而不是速度。另一方面，\hlblue{物体在任一给定时刻的运动性质，由它在该时刻速度的方向和大小决定（速度矢量永远沿物体轨迹的切线方向）。}

由于加速度和速度是不同的矢量，在一般情形下，作用力的方向与物体运动的方向可能不一致。因而，\hlblue{物体在给定时刻的运动性质，并不由该时刻作用在物体上的力唯一地决定。}

\std 这是一般情形。当然，作用力的方向与速度的方向也可能一致。

\tchr 当然可能。把一个物体举起来，然后小心地释放，使它不获得初速度：这时运动方向与重力方向一致。但如果你给物体一个水平方向的初速度，运动方向就不再与重力方向一致，物体将沿抛物线运动。虽然两种情形中物体都是在同一个力——重力——的作用下运动，运动的性质却不同。物理学家会说，这种差别源于初始条件（initial conditions）不同：运动开始时，第一种情形物体没有速度，第二种情形物体具有确定的水平速度。
\begin{figure}[!h]%
\centering
\includegraphics[width=0.4\textwidth]{figs/ch-04/fig-019a.pdf}
\caption[具有不同初速度的物体的轨迹。]{具有不同初速度的物体的轨迹。\label{fig-19}}
\end{figure}
\figr{fig-19} 画出了以不同方向的初速度抛出的各物体的轨迹；所有情形中，作用在物体上的都是同一个力——物体的重力。

\std 这是否意味着：物体在给定时刻的运动性质，不仅取决于该时刻作用在物体上的力，还取决于初始条件？

\tchr 正是如此。应当强调：初始条件反映着物体的既往历史。它们是过去存在过的力所产生的结果。这些力已不复存在，但其作用的结果仍显现着。从哲学观点看，这体现了过去与现在的联系，即因果性原理（principle of causality）。

请注意：倘若牛顿第二定律的公式里包含的是速度而不是加速度，过去与现在的这种联系就不会显露出来。那样的话，物体在给定时刻的速度（即该时刻运动的性质）就将完全由恰在此刻作用在物体上的力决定，过去对现在不会有任何影响。
\begin{figure}[!h]%
\centering
\includegraphics[width=0.4\textwidth]{figs/ch-04/fig-020a.pdf}
\caption[随初始条件不同，悬于绳上的小球可以有不同的合运动。]{随初始条件不同，悬于绳上的小球可以有不同的合运动。\label{fig-20}}
\end{figure}
我还想再举一个例子来说明上述道理，如\figr{fig-20} 所示：挂在绳上的小球受两个力的作用——重力和绳的张力。若把小球向平衡位置的一侧偏开然后释放，它就开始振动；若在垂直于偏离平面的方向上赋予小球一个确定的速度，它就开始沿圆周匀速运动。可见，随初始条件不同，小球或者在平面内振动（见\figr{fig-20}~\drkgry{(a)}），或者沿圆周匀速运动（见\figr{fig-20}~\drkgry{(b)}）。无论哪种情形，作用在小球上的都只有两个力：重力和绳的张力。

\std 我从来没有从这个角度考虑过牛顿定律。

\tchr 难怪有些学生在试图确定作用在物体上的力时，不去先弄清有哪些物体与该物体相互作用，而是从运动的性质出发进行推理。你可以回想起来，你自己就这么做过。正因如此，在画\figr{fig-08}~\drkgry{(c)} 与\figr{fig-08}~\drkgry{(d)} 时，你会觉得这两种情形中作用在物体上的力应该是不同的。实际上，两种情形中作用在物体上的都是两个力：重力和绳的张力。

\std 现在我明白了：\hlblue{同一组力可以引起性质不同的运动，因此，关于物体运动性质的资料不能作为确定作用在物体上的力的出发点。}

\tchr 你表述得非常准确。不过也不必走向极端。虽然同一组力可以引起不同种类的运动（如\figr{fig-20} 所示），但对不同的运动，各作用力之间的数值关系是不同的。这就是说，每一种运动对应的合力各不相同。譬如，物体沿圆周匀速运动时，合力应当是向心力；在平面内振动时，合力应当是回复力。由此可见，关于运动种类的资料虽然不能作为确定作用力的根据，却也远非多余。

顺着这个思路，让我们回到\figr{fig-20} 所示的例子。设已知绳的方向与竖直方向之间的夹角 $\alpha$，以及物体的重力 $P$。求下列两种情形中绳的张力 $T$：
\begin{enumerate*}[label=(\arabic*)]
\item 振动的物体处于极端位置时；
\item 物体沿圆周匀速运动时。
\end{enumerate*}

第一种情形，合力是回复力，方向垂直于绳。因此，把物体的重力 $P$ 分解为两个分量：一个沿合力方向，另一个垂直于合力（即沿绳的方向）。然后令垂直于合力的那些力——即沿绳方向作用的各力——相互平衡（见\figr{fig-21}~\drkgry{(a)}）。于是
\begin{equation*}%
T_{1} = P \cos \alpha
\end{equation*}
\begin{figure}%
\centering
\includegraphics[width=0.7\textwidth]{figs/ch-04/fig-021a.pdf}
\caption[把运动摆所受的力进行分解。]{把运动摆所受的力进行分解。\label{fig-21}}
\end{figure}
第二种情形，合力是向心力，方向水平。因此，应把绳的张力 $T_{2}$ 分解为一个竖直力和一个水平力，再令垂直于合力的那些力——即竖直方向的各力——相互平衡（\figr{fig-21}~\drkgry{(b)}）。

于是
\begin{equation*}%
T_{2} \cos \alpha = P \quad \text{或} \quad T_{2} = \frac{P}{\cos \alpha}
\end{equation*}
可见，了解物体运动的性质，对求绳的张力确实有用。
\pnote{第二种情形就是中学熟悉的圆锥摆（conical pendulum）：张力沿绳，其竖直分量与重力平衡（$T\cos\alpha=P$），水平分量提供向心力（$T\sin\alpha=mv^{2}/r$）；人教版必修第二册第六章“向心力”一节有“用圆锥摆粗略验证向心力表达式”的实验。第一种情形对应摆角最大处：此刻速度为零、向心加速度为零，故 $T=P\cos\alpha$。}

\std 如果我理解得都对，那么：知道了物体的相互作用，就能求出作用在其中某个物体上的力；知道了这些力和初始条件，就能预言物体运动的性质（任意时刻速度的大小与方向）；另一方面，知道了物体运动的种类，就能确定作用在它上面的各力之间的关系。我这样推理对吗？

\tchr 完全正确。不过我们继续。我想提一道与牛顿第二定律有关的比较简单的题。\hlred{两个质量分别为 $M$ 与 $m$ 的物体，被举到地板上方同一高度，然后同时释放。若空气对两物体的阻力相同，它们会同时到达地板吗？为简单起见，我们假定空气阻力是恒定的。}

\std 既然两物体所受的空气阻力相同，就可以把它忽略。于是两个物体同时到达地板。

\tchr 你错了。你无权忽略空气阻力。取质量为 $M$ 的物体来看，它受两个力作用：重力 $Mg$ 和空气阻力 $F$。合力是 $Mg-F$。由此求出加速度：
\begin{equation*}%
a = \dfrac{Mg -F}{M} = g - \dfrac{F}{M}
\end{equation*}
由此可见，质量较大的物体加速度较大，因而先到达地板。
\pnote{原书设空气阻力恒定，是为了让计算一目了然。实际的空气阻力随速度增大而增大，落体最终会达到稳定的收尾速度（terminal velocity）——跳伞员张伞后近似匀速下降就是这个道理。定性结论不变：同样的阻力对质量大的物体影响更小，一般仍是质量大者先落地。}

我要再次强调：计算物体的加速度时，必须计入作用在它上面的所有力，也就是说，必须求出合力。与此相关，“动力”（driving force）这一术语的使用应受批评。这个术语是不恰当的：把它用于某个（或某些）力，我们似乎就把这个（或这些）力在使物体获得加速度方面的作用单独挑了出来，仿佛其余的力不那么重要似的。这是完全错误的。\hlblue{物体的运动是作用在它上面的全部力无一例外的共同作用的结果}（当然，初始条件也应计入）。
\pnote{中文课堂与教辅同样常用“动力／阻力”这类按效果区分的说法（如“摩擦力充当动力”）。定性讨论方向时无妨，但列方程时必须无遗漏地计入全部作用力——“动力”并不是一种可以单独抽取出来的力，这正是本段要纠正的习惯。}

现在来看一个关于牛顿第三定律的例子。马开始拉马车；于是马和马车以某一加速度运动起来。按照牛顿第三定律，无论马以多大的力拉马车，马车都反过来以恰好同样大的力、沿相反方向拉马。既然如此，\hlred{为什么马和马车会向前加速运动？请解释。}

\std 这个问题我从来没想过，不过我看不出有什么矛盾。如果马作用在马车上的力被马车作用在马上的力所平衡，加速度就难以解释了。但这两个力不可能相互抵消，因为它们作用在不同的物体上：一个作用在马上，另一个作用在马车上。

\begin{figure}%[-3cm]
\centering
\includegraphics[width=0.5\textwidth]{figs/ch-04/fig-022a.pdf}
\caption[为什么马和马车向前加速运动？]{为什么马和马车向前加速运动？\label{fig-22}}
\end{figure}

\tchr 你的解释适用于马车没有套在马上的情形。那时马蹬离马车，马车朝一个方向运动，马朝另一个方向运动。我提出的情形完全不同：马套在马车上，两者连在一起，作为一个整体系统运动。你提到的马与马车之间的相互作用力，作用在同一系统的不同部分上；在系统整体的运动中，这两个力可以看作相互平衡的力。这样看来，你还没有回答我的问题。

\std 那我就不明白问题出在哪里了。也许这里的作用没有被反作用完全平衡？要知道，马可是活的生物啊。

\tchr 别让你的想象脱缰了。你不过遇到了一点困难，就准备牺牲力学的一条基本定律。要回答我的问题，完全不需要修改牛顿第三定律；相反，让我们恰恰以这条定律为讨论的基础。

按照第三定律，马与马车的相互作用不可能导致这一系统整体的运动（更确切地说，不可能使系统整体获得加速度）。既然如此，就必定还存在某种补充的相互作用。换句话说，除马和马车之外，问题中至少还须有一个物体参与。在给定情形中，这个物体就是地球。于是，我们要处理的相互作用就不是一种，而是三种：
\begin{enumerate}[label=(\arabic*),leftmargin=1cm]
\item 马与马车之间（把这个力记作 $f_{0}$）；
\item 马与地球之间（力 $F$），即马蹬地面的力；
\item 马车与地球之间（力 $f$），即马车与地面之间的摩擦力。
\end{enumerate}
所有物体都画在\figr{fig-22} 中：马、马车和地球；而且每个物体上都作用着两个力，它们是该物体与另外两个物体相互作用的结果。马—车系统的加速度由作用在该系统上全部力的合力引起；这样的力共有四个，其合力为 $F-f$。正是它使系统获得加速度。现在你看到了：这一加速度与马和马车之间的相互作用无关。
\pnote{马拉车问题是辨析牛顿第三定律的经典情境（人教版必修第一册第三章 3.3 节“牛顿第三定律”）：作用力与反作用力等大反向，但作用在不同物体上，本说不上“抵消”。本段再进一步：从“马＋车”整体看，这对相互作用力是内力，使整体加速的是地面给马的向前的力与地面给马车的向后的摩擦力之差 $F-f$。“内力不能改变系统整体的加速度”正是整体法列方程的依据（第 2 章关于整体法与隔离法的评注可对照）。}

\std 这么说来，地球表面不仅仅是某些事件发生的场所，而且是这些事件的积极参与者。

\tchr 你这个形象的说法很对。顺便说一句：如果把马和马车放到理想的光滑冰面上，从而排除该系统与地球之间的一切水平相互作用，那就根本不会有任何运动。应当强调：任何内部相互作用都不能使系统整体获得加速度，这件事只有靠外部作用才能做到（你不可能揪着自己的头发把自己提起来，也不可能拽着靴带把自己提起来）。这是牛顿第三定律的一个重要实际推论。

\end{dialogue}
\cleardoublepage
\thispagestyle{empty}


\thispagestyle{empty}
\vspace*{-2cm}
\begin{figure*}[!ht]
\centering
\includegraphics[width=0.8\textwidth]{figs/sec/tarasov-qasp-sec-03.pdf}
\end{figure*}
\begin{fullwidth}
{\textsf{\Large \hlred{力学掌握得好，习题就容易解。反过来同样成立：习题解得轻松，显然就说明力学掌握得好。因此，请通过尽可能多地解题来扩展你的力学知识。}}}
\end{fullwidth}

% NEWTERMS: frame of reference (reference system) | 参考系 | 人教版必修第一册第一章术语；本书首现于本章
% NEWTERMS: inertial frame of reference | 惯性参考系（惯性系） | 中学不作要求；大学普通物理/竞赛内容
% NEWTERMS: non-inertial frame of reference | 非惯性参考系（非惯性系） | 同上
% NEWTERMS: inertial force | 惯性力 | 中学不作要求；方向辨析见本章评注
% NEWTERMS: proportionality factor | 比例系数 | 牛顿第二定律 $B$ 系数语境；SI 中 $B=1$
% NEWTERMS: vector (quantity) | 矢量 | 基础术语，首现括注
% NEWTERMS: centripetal acceleration | 向心加速度 | 人教版必修第二册
% NEWTERMS: initial conditions | 初始条件 | 本章核心概念
% NEWTERMS: principle of causality | 因果性原理 | 哲学表述，照译
% NEWTERMS: driving force | 动力 | 按效果命名；原书批评其使用，评注对接中文“动力/阻力”说法
% NEWTERMS: extreme position | 极端位置 | 摆的最大偏角处
% NEWTERMS: conical pendulum | 圆锥摆 | 评注中给出；人教版必修第二册第六章
% NEWTERMS: terminal velocity | 收尾速度 | 评注中给出
